\section{The orientation guarantee}
\label{sec:orientation}

The solver uses an orientation to choose which boundary to truncate.
Choice \(Z\) lowers the upper bound and can exclude a subsolution;
choice \(Y\) raises the lower bound and can exclude a supersolution.
The alternative below applies when a fixed comparison has sufficiently
heavy support near the affected boundary: with high probability, the
procedure either chooses the other direction or returns a set on which
that support is concentrated enough for reweighting.
Section~\ref{sec:solving} shows that tightening ensures this hypothesis
whenever truncation threatens the comparison.

Fix an integer rank \(r\geq1\). Throughout this section, assume
Equation~\eqref{eq:comparison-guarantees} for every \(\Solve\) input of
rank less than \(r\), and consider \(\Orient\) inputs of rank \(r\).
The support parameters remain unchanged in an orientation's recursive
call, so that call retains rank \(r\). We prove its guarantee by
induction on the reference width. Calls from the center loop to
\(\Ask\), on the other hand, reduce rank and may use the assumed
comparison guarantees. As in Equation~\eqref{eq:global-parameters}, write
\[
  \varepsilon=2^{-J}.
\]
This error parameter is used only in the proof.

For a box \([a,b]\) with reference width \(B\), define the extreme
cores of a subsolution \(y\) and a supersolution \(z\) by
\begin{equation}
\label{eq:orientation-cores}
  I_y=\{i:y_i\geq a_i+B-hB\},
  \qquad
  I_z=\{i:z_i\leq b_i-B+hB\}.
\end{equation}
The inequalities \(B>0\) and \(h<1\) imply
\(I_y\subseteq S_y\) and \(I_z\subseteq S_z\).

\begin{proposition}[Orientation]
\label{prop:orientation}
Assume the lower-rank comparison guarantees stated at the beginning of
this section. On a fixed permitted input
\((a,b;B,\alpha,\beta,k,m)\) of rank \(r\), let
\((\chi,\mathcal H)\) be the output of \(\Orient\).
For every fixed admissible subsolution \(y\) of \([a,b]\) whose core
in Equation~\eqref{eq:orientation-cores} satisfies
\(\alpha(I_y)>\eta k\),
\begin{equation}
\label{ori:positive-alternative}
  \Pr\left(
     \chi=Y
     \ \text{or}\
     \exists H\in\mathcal H:
          \alpha(S_y\cap H)>\eta k
  \right)\geq 1-\frac1{16}.
\end{equation}
Separately, for every fixed admissible supersolution \(z\) of
\([a,b]\) with \(\beta(I_z)>\eta m\),
\begin{equation}
\label{ori:negative-alternative}
  \Pr\left(
     \chi=Z
     \ \text{or}\
     \exists H\in\mathcal H:
          \beta(S_z\cap H)>\eta m
  \right)\geq 1-\frac1{16}.
\end{equation}
On every tape, every set in \(\mathcal H\) has \(\mu\)-measure at
most \(1/2\). Each probability statement remains valid conditional
on a pre-call history that fixes its input and qualifying comparison,
when the call uses fresh random bits.
\end{proposition}

The two assertions are separate pointwise guarantees: they do not
require one event that works for all comparisons. Every returned set
\(H\) has \(\mu\)-measure at most \(1/2\), so a heavy support
intersection is exactly what Lemma~\ref{lem:reweighting} needs to
preserve the comparison while decreasing rank. We first prove the
proposition for one orientation trial, then amplify its alternative
for \(\OrientPlus\).

The proof follows the deviation of a fixed comparison from the centers
produced by the procedure. Reliable lower-rank comparisons make this
deviation nonincreasing. Whenever the projected support meets the
reduced budget, both tests qualify and force a definite decrease.
Many such times move the entire extreme core into a recorded
displacement set. If few times qualify, a fresh sampled center usually
has a heavy projected support; we will place that support in the
extreme core of the smaller recursive box. The next lemmas establish
these two routes before the width-induction proof.

\subsection{Whole-vector amplification}

\begin{lemma}[Amplified comparisons]
\label{lem:ask-amplification}
Let \(\Ask\) have an input of rank less than \(r\), and let \(p\) be
its returned vector. For each fixed admissible subsolution \(y\) of its
input box, and separately for each fixed admissible supersolution \(z\),
\[
  \Pr(p\not\geq y)\leq\varepsilon,
  \qquad
  \Pr(p\not\leq z)\leq\varepsilon.
\]
Both conclusions also hold conditional on any pre-call history that
determines the input and the indicated admissible comparison, provided
the call uses fresh randomness.
\end{lemma}

\begin{proof}
Fix the input and \(y\). Count a trial as successful only when its
entire output dominates \(y\). If a strict majority of the \(J\)
trials succeeds, then at each coordinate a strict majority of the
entries is at least the corresponding coordinate of \(y\). The
coordinatewise median therefore dominates \(y\).
Consequently its failure requires at least
\(t=(J+1)/2\) unsuccessful trials. The lower-rank hypothesis and
independence bound the probability that any specified \(t\) trials
all fail by \(16^{-t}\). Taking a union bound over at most \(2^J\)
choices of such trials gives
\[
  \Pr(p\not\geq y)
  \leq 2^J16^{-t}
  \leq 2^J16^{-J/2}
  =2^{-J}.
\]
For a fixed \(z\), count a trial as successful when its entire output
is at most \(z\). A strict majority of these trials places the median
below \(z\) in every coordinate, and the same count proves its bound.

To apply either argument conditionally, first fix the history preceding
the call. It specifies a deterministic input and a deterministic
comparison. Whenever that comparison is admissible for the specified
input, the fresh independent trial bits give exactly the experiment
just analyzed. Averaging these conditional bounds gives the
unconditional assertion. The proof uses whole-vector trial events;
it requires neither independence among coordinates nor an event valid
simultaneously for different comparison vectors.
\end{proof}

\subsection{Geometry of the centers and projections}

The loop tests comparisons in boxes of radius \(D\); a recursive
orientation call uses the larger radius \(E_0\). Both projections
must be admissible, and the support tested at radius \(D\) must lie
in the extreme core of the radius-\(E_0\) box. The next estimates
supply these facts: the deviations are large enough for both projections, and the
boundary layer of the larger box is at least \(2D\) wide. Only a
nonempty core is needed here, and the estimates hold on every tape.

\begin{lemma}[Center estimates]
\label{ori:center-estimates}
Let \(B>128s\), and let \(q_0,\ldots,q_s\) be the centers in
Step~\ref{alg:orient-loop} of \(\Orient\). Put
\[
  D=\left\lfloor\frac{B}{32s}\right\rfloor,
  \qquad
  E_0=\left\lfloor\frac B8\right\rfloor.
\]
Then \(q_j\in[a,b]\) and
\(\lvert q_{j+1,i}-q_{j,i}\rvert\leq D\) for every applicable
\(i,j\), on every random tape. Moreover,
\begin{equation}
\label{ori:radius-estimates}
  D\geq\frac{B}{64s}>0,
  \qquad
  sD\leq E_0,
  \qquad
  2D\leq h(2E_0).
\end{equation}
For a subsolution \(y\) with \(I_y\ne\varnothing\), set
\(d_j=\max_i(y_i-q_{j,i})\). Then
\begin{align}
  d_j&\geq B/2-hB-sD-1\geq 2E_0
       &&(0\leq j\leq s), \label{ori:positive-deviation}\\
  d_0-(y_i-q_{0,i})&\leq hB+1
       &&(i\in I_y). \label{ori:positive-deficit}
\end{align}
For a supersolution \(z\) with \(I_z\ne\varnothing\), set
\(e_j=\max_i(q_{j,i}-z_i)\). Then
\begin{align}
  e_j&\geq B/2-hB-sD-1\geq 2E_0
       &&(0\leq j\leq s), \label{ori:negative-deviation}\\
  e_0-(q_{0,i}-z_i)&\leq hB+1
       &&(i\in I_z). \label{ori:negative-deficit}
\end{align}
The coordinate widths may be unequal and may be zero.
\end{lemma}

\begin{proof}
The output-type assertion of Proposition~\ref{prop:termination} puts
both answers at time \(j\), and therefore their floored average,
in \([q_j-D,q_j+D]\). Since \(q_j\in[a,b]\), clipping that average
to \([a,b]\) cannot move a coordinate farther from \(q_j\).
Starting with \(q_0\in[a,b]\), induction proves containment and
the step bound without any comparison-success assumption.

For the radius estimates, \(B/(32s)>4\), so its floor is at least
half its value. Also,
\[
  sD\leq \frac{B}{32}
       \leq \frac B8-1
       \leq E_0.
\]
Here \(B>128s\) and the global bound \(s\geq16\) justify the
middle inequality. Dividing the resulting inequality \(sD\leq E_0\)
by \(s\) gives the last inequality in
Equation~\eqref{ori:radius-estimates}.

To compare a core coordinate with the initial center, write
\(\ell_i=b_i-a_i\). Then \(0\leq\ell_i\leq B\) and
\[
  q_{0,i}=a_i+\left\lfloor\frac{\ell_i}{2}\right\rfloor.
\]
On the subsolution side, \(y_i\leq b_i\) implies
\[
  d_0\leq\max_i\left\lceil\frac{\ell_i}{2}\right\rceil
      \leq\left\lceil\frac B2\right\rceil.
\]
A coordinate of the high core, however, satisfies
\[
  y_i-q_{0,i}
  \geq B-hB-\left\lfloor\frac{\ell_i}{2}\right\rfloor
  \geq \frac B2-hB.
\]
The difference between these two bounds is at most \(hB+1/2\),
which proves Equation~\eqref{ori:positive-deficit}. During the loop
each coordinate moves by at most \(sD\) in total. Keeping any one
high-core coordinate in the maximum defining \(d_j\) therefore gives
\(d_j\geq B/2-hB-sD\), and in particular the first inequality of
Equation~\eqref{ori:positive-deviation}.

For the supersolution, \(z_i\geq a_i\) instead gives
\[
  e_0\leq\max_i\left\lfloor\frac{\ell_i}{2}\right\rfloor
      \leq\left\lfloor\frac B2\right\rfloor.
\]
At a coordinate in the low core,
\[
  q_{0,i}-z_i
  \geq B-hB-\left\lceil\frac{\ell_i}{2}\right\rceil
  \geq \frac B2-hB-\frac12.
\]
This proves the initial-deficit bound
\eqref{ori:negative-deficit}. Subtracting the possible movement
\(sD\) proves the first inequality in
Equation~\eqref{ori:negative-deviation}.

It remains to place both deviations above the child-box width. Using
\(sD\leq B/32\) and \(2E_0\leq B/4\), we obtain
\[
  \left(\frac B2-hB-sD-1\right)-2E_0
  \geq B\left(\frac7{32}-\frac1s\right)-1
  >0.
\]
The final inequality follows from \(s\geq16\) and \(B>128s\).
It completes both deviation estimates.
\end{proof}

These estimates let us use the centered projections from
Lemma~\ref{lem:center-projections} at both radii needed by the
algorithm: \(D\) in the loop and \(E_0\) in the recursive call.
Whenever the corresponding core is nonempty, denote these projections by
\begin{equation}
\label{ori:projected-comparisons}
  y^{j,A}=\max(y-d_j+A,q_j-A),
  \qquad
  z^{j,A}=\min(z+e_j-A,q_j+A).
\end{equation}
For \(A\in\{D,E_0\}\), they are comparisons in
\([q_j-A,q_j+A]\), with exact supports
\begin{equation}
\label{ori:projected-supports}
  S_{y^{j,A}}=\{i:y_i-q_{j,i}>d_j-2A\}\subseteq S_y,
  \qquad
  S_{z^{j,A}}=\{i:q_{j,i}-z_i>e_j-2A\}\subseteq S_z.
\end{equation}
Thus each projection keeps the parent's admissibility under unchanged
support weights and budget. The argument uses \(q_j\in[a,b]\) and
the deviation lower bounds; it does not require the centered child
box to be contained in \([a,b]\).

\subsection{Reduced budgets force a decrease in deviation}

Consider one of these fixed admissible comparisons. At each time its
projection qualifies for the \(\Ask\) call that retains its own
budget. Sometimes it also qualifies for the call that reduces that
budget. We call those times \emph{light}, referring to the weight of
the projected support. For \(y\), the condition is
\begin{equation}
\label{ori:bad-positive}
  \alpha\bigl(\{i:y_i-q_{j,i}>d_j-2D\}\bigr)\leq\eta k.
\end{equation}
For \(z\), it is
\begin{equation}
\label{ori:bad-negative}
  \beta\bigl(\{i:q_{j,i}-z_i>e_j-2D\}\bigr)\leq\eta m.
\end{equation}
These conditions are defined separately for the two comparisons.
Each is determined by the history before the two calls at time \(j\).
We next show why a light time forces progress for its fixed comparison.

\begin{lemma}[Deviation drops]
\label{ori:deviation-drops}
For each fixed admissible subsolution \(y\) with nonempty high core,
with probability at least \(1-2s\varepsilon\) every loop step satisfies
\[
  d_{j+1}\leq d_j,
  \qquad
  d_{j+1}\leq d_j-D\quad\hbox{at every light time for \(y\)}.
\]
Separately, for each fixed admissible supersolution \(z\) with nonempty
low core, with probability at least \(1-2s\varepsilon\) every step satisfies
\[
  e_{j+1}\leq e_j,
  \qquad
  e_{j+1}\leq e_j-D\quad\hbox{at every light time for \(z\)}.
\]
\end{lemma}

\begin{proof}
We first work with \(y\) and condition on the history before time
\(j\). The projection \(y^{j,D}\) is now fixed and is admissible
for the second call, which retains the subsolution budget \(k\) and
reduces \(m\). This call has rank less than \(r\), so
Lemma~\ref{lem:ask-amplification} gives, with conditional error at
most \(\varepsilon\),
\[
  r_2\geq y^{j,D}\geq y-d_j+D.
\]
The first answer satisfies, without a correctness assumption,
\[
  r_1\geq q_j-D\geq y-d_j-D.
\]
Averaging gives the integer lower bound \(y-d_j\), which is preserved
by flooring. Since \(b\geq y\geq y-d_j\), upper clipping preserves
it as well; lower clipping can only help. Thus the next center obeys
\(q_{j+1}\geq y-d_j\), equivalently \(d_{j+1}\leq d_j\).

At a light time, the exact support formula
\eqref{ori:projected-supports} and the qualification
\eqref{ori:bad-positive} also make \(y^{j,D}\) admissible for the
first call, whose budget is \(\eta k\). Allowing one more conditional
error of at most \(\varepsilon\), both answers are at least
\(y-d_j+D\). Their floored average retains this integer lower bound.
Furthermore, \(d_j\geq2E_0\geq2D\), so
\(y-d_j+D\leq y\leq b\); clipping also retains it. This yields
\(q_{j+1}\geq y-d_j+D\), or \(d_{j+1}\leq d_j-D\).

For the separate supersolution argument, condition on the same kind
of pre-step history. The first call retains the full supersolution
budget, so its comparison with \(z^{j,D}\) gives, except with
conditional probability \(\varepsilon\),
\[
  r_1\leq z^{j,D}\leq z+e_j-D.
\]
The second answer always satisfies
\[
  r_2\leq q_j+D\leq z+e_j+D.
\]
Their average is at most \(z+e_j\), and flooring preserves an upper
bound. Lower clipping preserves it because
\(a\leq z\leq z+e_j\); upper clipping can only help. We obtain
\(q_{j+1}\leq z+e_j\) and therefore \(e_{j+1}\leq e_j\).

At a light time for \(z\), Equation~\eqref{ori:bad-negative} lets
us request the comparison with \(z^{j,D}\) from the second call too.
With one additional conditional error of at most \(\varepsilon\),
both answers are at most \(z+e_j-D\). Flooring respects this bound,
and \(e_j\geq2D\) gives \(z+e_j-D\geq z\geq a\), so lower
clipping respects it. Hence \(e_{j+1}\leq e_j-D\).

For either chosen side, at most two amplified comparisons are requested
per step. The comparison vector and any reduced-budget qualification
are determined before the appropriate fresh call. Applying the
conditional assertion of Lemma~\ref{lem:ask-amplification} and summing
over the \(s\) steps bounds the chance of any required failure by
\(2s\varepsilon\).
\end{proof}

The decrease in maximum deviation translates into displacement of every
coordinate in the extreme core, because the initial deficits on that
core are bounded. Recall the algorithm's two displacement sets:
\[
  P=\{i:q_{s,i}-q_{0,i}>2hB\},
  \qquad
  N=\{i:q_{s,i}-q_{0,i}<-2hB\}.
\]

\begin{lemma}[Core displacement]
\label{ori:core-displacement}
On the subsolution event of Lemma~\ref{ori:deviation-drops},
\(I_y\cap N=\varnothing\); if there are more than \(256\) light
times for \(y\), then \(I_y\subseteq P\).
On the separate supersolution event,
\(I_z\cap P=\varnothing\); more than \(256\) light times for \(z\)
imply \(I_z\subseteq N\).
\end{lemma}

\begin{proof}
For \(i\in I_y\), compare the definition of \(d_s\) with the
initial-deficit estimate \eqref{ori:positive-deficit}:
\begin{align*}
  q_{s,i}-q_{0,i}
  &\geq y_i-d_s-q_{0,i}\\
  &=(d_0-d_s)-\bigl(d_0-(y_i-q_{0,i})\bigr)\\
  &\geq(d_0-d_s)-(hB+1).
\end{align*}
On the stipulated event, \(d_0-d_s\geq0\). Since \(hB=B/s>128\),
the last expression is greater than \(-2hB\), excluding \(N\).
More than \(256\) light times means at least \(257\) drops, and gives
\[
  d_0-d_s\geq257D\geq\frac{257}{64}hB.
\]
Substituting this stronger bound yields
\[
  q_{s,i}-q_{0,i}-2hB
  \geq \left(\frac{257}{64}-3\right)hB-1
  =\frac{65}{64}hB-1
  >0.
\]
The strict inequality is exactly what is needed for \(i\in P\).

For \(i\in I_z\), use \eqref{ori:negative-deficit} to get
\begin{align*}
  q_{s,i}-q_{0,i}
  &\leq z_i+e_s-q_{0,i}\\
  &=\bigl(e_0-(q_{0,i}-z_i)\bigr)-(e_0-e_s)\\
  &\leq(hB+1)-(e_0-e_s).
\end{align*}
Nonincrease of \(e_j\) makes this strictly less than \(2hB\),
excluding \(P\). With at least \(257\) light times,
\(e_0-e_s\geq257D\), so the same arithmetic gives
\[
  q_{s,i}-q_{0,i}+2hB
  \leq 1-\frac{65}{64}hB<0.
\]
Thus \(i\in N\), as asserted.
\end{proof}

\subsection{Proof of the orientation alternative}

The displacement lemma handles every early return. When the procedure
recurses, it also handles many light times by placing the heavy core
in an appended set. In the remaining case, fresh sampling usually
selects a projection whose support is heavy enough to invoke the
same alternative at the smaller width.

\begin{proof}[Proof of Proposition~\ref{prop:orientation}]
We prove a quantitative version of each pointwise statement by width
induction. Define the error allowance per nonterminal level by
\[
  \gamma=\frac{256}{s}+2s\varepsilon,
\]
and define the number of such levels recursively by
\[
  \lambda(B)=
  \begin{cases}
    0,&B\leq128s,\\
    1+\lambda\bigl(2\lfloor B/8\rfloor\bigr),&B>128s.
  \end{cases}
\]
A recursive width is at most \(B/4\), so \(\lambda(B)\leq g+1\).
We will bound the failure probability for each fixed qualifying
comparison by \(\lambda(B)\gamma\).

\paragraph{Exact orientation at small width.}
When \(B\leq128s\), the procedure computes the box solution \(p\)
and uses
\[
  P=\{i:p_i\geq a_i+B/2\}.
\]
Lemma~\ref{lem:box-comparison} gives \(y\leq p\leq z\). For a high-core
coordinate,
\(p_i\geq y_i\geq a_i+(1-h)B>a_i+B/2\); hence \(I_y\subseteq P\).
For a low-core coordinate, the reference-width bound gives
\[
  p_i\leq z_i\leq b_i-B+hB\leq a_i+hB<a_i+B/2,
\]
so \(I_z\subseteq V\setminus P\).
If \(\mu(P)>1/2\), the returned choice \(Y\) protects the
subsolution, and the returned set \(V\setminus P\) contains the low
core. Otherwise the choice \(Z\) protects the supersolution, and
the returned set \(P\) contains the high core. Core containment in
support transfers the assumed strict weight bound to the indicated
exceptional set. In either branch that set has measure at most
\(1/2\). Both alternatives therefore hold with zero error.

\paragraph{A fixed subsolution: returns explained by motion.}
Now take \(B>128s\) and fix \(y\) with \(\alpha(I_y)>\eta k\).
The core is nonempty. Charge an error of at most \(2s\varepsilon\)
for the event that the deviation conclusions fail, and first work on
the complementary event from Lemma~\ref{ori:deviation-drops}.
Lemma~\ref{ori:core-displacement} is then available for this \(y\).

If \(\mu(P)>1/2\), the procedure returns \(Y\), proving the
alternative directly. If it returns instead at the second test,
\(\mu(N)>1/2\), it supplies \(V\setminus N\). Since
\(I_y\cap N=\varnothing\),
\[
  \alpha(S_y\cap(V\setminus N))
  \geq\alpha(I_y)>\eta k.
\]
This exceptional set proves the alternative in the second early-return
case as well.

Suppose neither test returns. Both \(P\) and \(N\) have measure
at most \(1/2\) and will be appended to the child's list.
If there are more than \(256\) light times for \(y\), the displacement
lemma gives \(I_y\subseteq P\). Thus
\(\alpha(S_y\cap P)>\eta k\), and the appended set already proves
the parent assertion, whatever happens in the recursive call.

\paragraph{A fixed subsolution: transfer at the chosen iterate.}
It remains to handle a completed loop with at most \(256\) light
times. Fix its entire history. The index \(j\) is chosen only now,
using fresh bits, so its conditional distribution is uniform on
\(\{0,\ldots,s-1\}\). Its probability of being light is at most
\(256/s\). This remains true on the histories where all required
loop comparisons succeeded, since that event is determined before
sampling.

For a sampled index that is not light, write
\[
  T_y=\{i:y_i-q_{j,i}>d_j-2D\},
  \qquad
  B'=2E_0,
  \qquad
  a'=q_j-E_0,\quad b'=q_j+E_0.
\]
Then \(\alpha(T_y)>\eta k\). In the child box use the comparison
\[
  y'=y^{j,E_0}=\max(y-d_j+E_0,q_j-E_0).
\]
The projection formula \eqref{ori:projected-supports} makes \(y'\)
admissible and gives \(S_{y'}\subseteq S_y\). For \(i\in T_y\),
\begin{align*}
  y'_i
  &\geq y_i-d_j+E_0\\
  &>q_{j,i}+E_0-2D\\
  &\geq q_{j,i}+E_0-hB'
   =a'_i+B'-hB'.
\end{align*}
The last inequality uses \(2D\leq hB'\) from
Equation~\eqref{ori:radius-estimates}. It places all of \(T_y\)
in the child's high core, so \(\alpha(I_{y'})>\eta k\).
In particular, the strict inequality in the preceding line is
compatible with the inclusive threshold defining that core.

The completed loop and the sampled index together determine the child
input and \(y'\). Fix this information before drawing the child's
fresh bits. The width induction applies conditionally, with failure
probability at most \(\lambda(B')\gamma\).
On success the child either returns \(Y\), which the parent copies,
or supplies a set \(H\) satisfying
\[
  \alpha(S_{y'}\cap H)>\eta k.
\]
The parent retains \(H\). Since \(S_{y'}\subseteq S_y\) and the
support weights and budgets have not changed,
\(\alpha(S_y\cap H)>\eta k\) follows with the same strict threshold.

The motion argument has already settled all successful-loop histories
with more than \(256\) light times. On the remaining successful-loop
histories, the conditional sampling error is at most \(256/s\), and
the conditional child error is at most \(\lambda(B')\gamma\).
Averaging these conditional bounds and adding the loop error gives
a total failure probability for the fixed \(y\) of at most
\[
  2s\varepsilon+\frac{256}{s}+\lambda(B')\gamma
  =\lambda(B)\gamma.
\]

\paragraph{A fixed supersolution: returns explained by motion.}
Fix now \(z\) with \(\beta(I_z)>\eta m\), for a separate analysis.
Its core is nonempty. Outside an error event of probability at most
\(2s\varepsilon\), the corresponding deviation and displacement
conclusions hold; in particular \(I_z\cap P=\varnothing\).
A first-test return of \(Y\) supplies \(V\setminus P\), which
contains \(I_z\) and therefore satisfies
\(\beta(S_z\cap(V\setminus P))>\eta m\).
A second-test return chooses \(Z\) and directly proves the alternative.

If neither return occurs, both displacement sets have measure at most
\(1/2\). More than \(256\) light times now give \(I_z\subseteq N\).
The appended set \(N\) then satisfies
\(\beta(S_z\cap N)>\eta m\), independently of the child's outcome.

\paragraph{A fixed supersolution: transfer at the chosen iterate.}
With at most \(256\) light times, condition on the completed loop as
before. The fresh uniform index is light with probability at most
\(256/s\). For an index that is not light, put
\[
  T_z=\{i:q_{j,i}-z_i>e_j-2D\},
  \qquad
  z'=z^{j,E_0}=\min(z+e_j-E_0,q_j+E_0).
\]
Here \(\beta(T_z)>\eta m\). The projection \(z'\) is an admissible
supersolution in \([a',b']=[q_j-E_0,q_j+E_0]\), with reference width
\(B'=2E_0\), and \(S_{z'}\subseteq S_z\). For \(i\in T_z\),
\begin{align*}
  z'_i
  &\leq z_i+e_j-E_0\\
  &<q_{j,i}-E_0+2D\\
  &\leq q_{j,i}-E_0+hB'
   =b'_i-B'+hB'.
\end{align*}
Thus \(T_z\) lies in the child's low core, giving
\(\beta(I_{z'})>\eta m\). Once the loop history and \(j\) have
been fixed, this is a qualified comparison determined before the
recursive randomness. The induction hypothesis therefore applies.
A successful child returns \(Z\), copied by the parent, or a retained
set \(H\) with \(\beta(S_{z'}\cap H)>\eta m\). Support containment
then gives \(\beta(S_z\cap H)>\eta m\). Adding the loop, sampling,
and child errors proves the same bound \(\lambda(B)\gamma\) for
this fixed supersolution.

\paragraph{Measures, total error, and conditioning.}
The measure assertion holds independently of either comparison event.
Every early return uses a set of measure at most \(1/2\), or the
complement of a set of measure greater than \(1/2\). A recursive
return appends only \(P,N\) of measure at most \(1/2\).
Its child retains \(\alpha,\beta,k,m\), and hence the same measure
\(\mu\); the measure bounds for its sets transfer unchanged.
This proves the assertion by width induction, also in agreement with
Proposition~\ref{prop:termination}.

The accumulated error is bounded by
\[
  \lambda(B)\gamma
  \leq(g+1)\left(\frac{256}{s}+2s\,2^{-J}\right)
  \leq\frac1{512}+J\,2^{-J}
  <\frac1{16}.
\]
Indeed, \(s\geq2^{17}(g+1)\) controls the sampling term, while
\(2s(g+1)\leq s^2\leq J\) controls the loop term. Finally,
\(J\geq16\) gives \(J2^{-J}\leq16\,2^{-16}=1/4096\).
This closes both width inductions.

At each probabilistic step we fixed the relevant input and one
qualified comparison before exposing fresh bits. The same sequence
of arguments therefore applies after fixing any history before the
initial call that determines its input and qualifying comparison.
This proves the conditional assertions as well.
\end{proof}

\subsection{Amplifying the orientation alternative}

A trial can succeed either through its choice or through its list.
The majority operation applies to the choices, while concatenation
preserves the possible certificates from every trial. The following
argument uses precisely this output rule, including the lists of
minority trials.

\begin{lemma}[Amplified orientation]
\label{lem:orient-plus}
Under the lower-rank assumption of this section, the output of
\(\OrientPlus\) satisfies Equation~\eqref{ori:positive-alternative}
for each fixed qualifying subsolution and, separately,
Equation~\eqref{ori:negative-alternative} for each fixed qualifying
supersolution, with failure probability at most \(\varepsilon\)
in each case. The same bounds hold conditional on a history that
fixes the input and comparison before the call's fresh randomness.
On every tape its list has length at most \(J(2g+2)\), and each listed
set has measure at most \(1/2\).
\end{lemma}

\begin{proof}
Fix an input and a qualifying \(y\). Failure of the amplified
alternative means that the majority choice is \(Z\) and that no
set anywhere in the concatenated list has
\(\alpha(S_y\cap H)>\eta k\).
Every trial voting \(Z\) consequently failed its own alternative:
it neither chose \(Y\) nor returned a qualifying set. There are at
least \((J+1)/2\) such trials. By
Proposition~\ref{prop:orientation}, each trial has failure probability
at most \(1/16\). Independence and the same subset count used for
median amplification give
\[
  \Pr(\text{amplified failure for }y)
  \leq 2^J(1/16)^{J/2}
  =\varepsilon.
\]
For a fixed qualifying \(z\), failure instead means a majority
choice \(Y\) and no set with heavy intersection with \(S_z\).
Every \(Y\)-voting trial then failed the supersolution alternative.
There are again at least \((J+1)/2\) such trials, giving the same
probability bound.

Because all lists are retained, including those of minority trials,
the absence of a qualifying set from the amplified list implies its
absence from every trial's list. Each retained entry has measure at
most \(1/2\) by Proposition~\ref{prop:orientation}. There are at
most \(2g+2\) entries per trial by
Proposition~\ref{prop:termination}, so concatenating the \(J\) lists
has the asserted length even when entries repeat.
Lastly, fixing a pre-call history fixes the trial inputs and the
chosen comparison. The trial bits remain fresh and independent,
so the same amplification proof gives the conditional bound.
\end{proof}
